%%%
%
% $Autor: Wings $
% $Datum: 2021-05-14 $
% $Pfad:  $
% $Dateiname: SBOM
% $Version: 4620 $
%
% !TeX spellcheck = de_DE/GB
% !TeX program = pdflatex
% !BIB program = biber/bibtex
% !TeX encoding = utf8
%
%%%


\chapter{SBOM}

\section{Software Requirements}
The project depends on publicly available, free-of-cost software libraries and tools.

\subsection{Python Dependencies}
The following Python libraries are used in the project:

\begin{itemize}
    \item \textbf{logging} – Used for recording events and debugging information.
    \item \textbf{pandas} – A data analysis and manipulation library.
    \item \textbf{pykalman} – Implements the Kalman filter.
    \item \textbf{tensorflow} – A machine learning framework used in model training, inference, and optimization.
\end{itemize}

All these Python libraries are open-source and can be installed using the Python package manager \texttt{pip}.

\subsection{Arduino IDE Dependencies}
The Arduino scripts require the following libraries, which are also freely available:

\begin{itemize}
    \item \texttt{TensorFlowLite.h} – TensorFlow Lite for microcontrollers, used for on-device machine learning.
    \item \texttt{tensorflow/lite/micro/all\_ops\_resolver.h} – Registers all available TensorFlow Lite operators.
    \item \texttt{tensorflow/lite/micro/micro\_interpreter.h} – Runs inference on the microcontroller.
    \item \texttt{tensorflow/lite/schema/schema\_generated.h} – Defines the structure of TensorFlow Lite models.
    \item \texttt{tensorflow/lite/version.h} – Contains TensorFlow Lite versioning information.
    \item \texttt{Arduino\_LSM9DS1.h} – Library for the LSM9DS1 sensor.
\end{itemize}

These Arduino libraries are open-source and can be installed via the Arduino Library Manager.

\section{Availability and Licensing}
All software used in this project is **free of cost** and available to the public under open-source licenses. The Python libraries can be installed from the Python Package Index (PyPI), and the Arduino libraries are available through the Arduino Library Manager.

\section{Installation Guide}
To install the required Python dependencies, run:

\begin{verbatim}
pip install -r requirements.txt
\end{verbatim}

For Arduino, install the required libraries via Sketch > Include Library > Manage Libraries in the Arduino IDE.

\section{Usage}
Ensure all dependencies are installed before running the project. The scripts rely on the listed packages for data handling, filtering, and machine learning functionalities.

For the complete list of requirements, refer to the file below:

\href{run:../Code/requirements.txt}{Open requirements.txt}
